% GENERATED FILE. Edit canonical lessons, then run npm run generate:books.
% canonical-source-hash: fnv1a64:81c69340c1ca4a9d
% canonical-lessons: ES-C45-os

\chapter{Os --- You All, in Spain}
\label{ch:pronoun-os}
\hlchaptermodality{90}

\begin{chapteropening}
\textbf{By the end of this chapter:} \emph{I can use os for a group in Spain and los or las everywhere else.}

Complete ES-C45-os: say os quiero and its non-peninsular equivalent, and name the older word os is worn down from.
\end{chapteropening}

\section[os]{os — \textquotedblleft{}you all,\textquotedblright{} in Spain}

\label{lesson:ES-C45-os}



\begin{quote}
Which word for \textquotedblleft{}you all\textquotedblright{} is used in Spain, and which everywhere else? (\emph{Vosotros} in Spain; \emph{ustedes} everywhere.) That split is about to appear one last time.
\end{quote}

\begin{grammarlens}[title={Grammar Lens: the object form of vosotros}]
\begin{quote}
\textbf{os} — \emph{you all}, as the one it happens to.
\end{quote}

\begin{quote}
\emph{\textbf{Os} quiero.} — I love you all.
\end{quote}

Two letters, the shortest word in the set. It pairs with \emph{vosotros}, so it lives exactly where \emph{vosotros} lives: Spain, and nowhere else in daily speech.

Across the Atlantic, \textquotedblleft{}you all\textquotedblright{} as an object is \textbf{los} or \textbf{las} — the same forms you learned last chapter, because \emph{ustedes} is grammatically a \emph{they}-word:

\begin{quote}
\emph{Os quiero.} — I love you all. \textbf{(Spain)} \emph{Los quiero.} — I love you all. \textbf{(everywhere else)}
\end{quote}

Neither is a mistake for the other. Several hundred million speakers will never say \emph{os} in their lives, and forty million say it every day.
\end{grammarlens}

\begin{cousinweb}
\emph{Os} was \textbf{vos}, and the \emph{v} simply fell off in the object slot — a word this small, said this often, and never stressed, wears faster than the rest of the language.

That same \emph{vos} is still visible in two places you have met. It is the front of \textbf{vosotros} (\emph{vos} + \emph{otros}, \textquotedblleft{}you others,\textquotedblright{} built exactly like \emph{nosotros}). And it survives \textbf{whole} as a subject pronoun in Argentina and much of Central America, where \emph{vos} replaces \emph{tú}.

So one Latin word left three descendants alive in three different conditions: worn to \emph{os} in Spain, glued into \emph{vosotros} in Spain, and standing intact as \emph{vos} across the Río de la Plata. Same word, three fates, decided by geography.
\end{cousinweb}

\subsection*{Your turn}
Say these aloud:

\begin{itemize}
\raggedright
  \item \textquotedblleft{}Os quiero\textquotedblright{} then \textquotedblleft{}Los quiero\textquotedblright{} — the same sentence, two continents
  \item \textquotedblleft{}vos\textquotedblright{}, \textquotedblleft{}os\textquotedblright{}, \textquotedblleft{}vosotros\textquotedblright{} and hear one word in three states
\end{itemize}

\begin{tcolorbox}[breakable,colback=teal!4,colframe=teal!35!black,title={Before you move on}]
\textquotedblleft{}I love you all,\textquotedblright{} in Spain? (\emph{\textbf{Os quiero}}.) Elsewhere? (\emph{\textbf{Los quiero}}.) What is \emph{os} worn down from? (\emph{\textbf{Vos}}.) Next: all eight of these in one place.
\end{tcolorbox}
